\documentclass[a4paper]{article}

\usepackage{graphicx}
\usepackage{amsmath,amssymb}
\usepackage{minted}
\usepackage{multirow}
\usepackage{float}
\usepackage{algorithm}
\usepackage{algpseudocode}
\usepackage{todonotes}
\usepackage{forest}
\usepackage{xstring}
\usepackage{xcolor}
\usepackage[most]{tcolorbox}
\usepackage{xparse}
\usepackage{natbib}
\usepackage{adjustbox}

\usetikzlibrary{graphs,graphdrawing}
\usegdlibrary{layered}

% for external graphviz
\input{/nd_language_ai_project/src/nd_language_ai/formal/computer/programming/tex/latex/code_clips/external_graphviz.tex}

% for tags
\input{/nd_language_ai_project/src/nd_language_ai/formal/computer/programming/tex/latex/part/control_sequence/kind/semtag/semtag.tex}

% for links
\input{/nd_language_ai_project/src/nd_language_ai/formal/computer/programming/tex/latex/code_clips/seehere.tex}

% for nested lists and sections
\input{/nd_language_ai_project/src/nd_language_ai/formal/computer/programming/tex/latex/code_clips/deep_structure.tex}

\title{}
\date{}
\author{Mohammad Rahmani}

\begin{document}
   \maketitle

   Gestern feierten wir den 30. Geburtstag meiner Cousine Katrin. Die meisten Gäste waren für 15 Uhr eingeladen, aber ich war schon früher da, um bei den Vorbereitungen zu helfen. Ich deckte den Tisch, stellte die Kerzen auf und schmückte die ganze Wohnung mit Luftballons. Und dann klingelte es auch schon an der Tür.

   Katrin erhielt viele Glückwunschkarten, Blumensträuße und tolle Geschenke. Und jeder Gast versicherte ihr: „Du siehst so jung aus, dein Alter sieht man dir gar nicht an!“ Als alle Gäste da waren, setzten wir uns, um Katrins selbstgebackene Kuchen und Torten zu probieren. Da neben der Familie auch viele Freunde gekommen waren, wurde es ein bisschen eng in der Wohnung, aber jeder fand einen Platz.

   Gegen 19 Uhr war es Zeit fürs Abendessen. Katrin und ich hatten in der Küche ein Buffet vorbereitet, an dem sich jeder bedienen konnte. Ganz typisch für Deutschland sind Kartoffelsalat und Würstchen, die durften natürlich auch auf unserem Buffet nicht fehlen.

   Als ich mit meinem Teller aus der Küche kam, wunderte ich mich, warum im Flur Geschenkband und zerrissenes Geschenkpapier lagen. Ich folgte der Spur und entdeckte die Übeltäter: Katrins Katzen hatten ein kleines Geschenk geöffnet. Und während eine Katze mit der Schleife spielte, wollte die andere gerade das Geschenk fressen – ein Stück Seife. Schnell nahm ich die Seife weg, die war allerdings schon ein wenig angefressen. Ich hoffe nur, dass keine Seifenblasen entstehen, wenn die Katze das nächste Mal Wasser trinkt!
   \bibliographystyle{plainnat}
   \bibliography{/home/donkarlo/Dropbox/repo/nd_shared_research/bibliography/bibliography.bib}
\end{document}
